\documentclass{exam}
\usepackage{../commonheader}

\discnumber{5}
\title{Recursion, Tree Recursion}
\date{September 21, 2026 - September 25, 2026}

\begin{document}
\maketitle

\begin{meta}
    \textbf{Example Timeline}
    \begin{itemize}
        \item Intros \& Logistics [5 min]
        \item Recursion Basics: base cases, recursive cases, and the recursive leap of faith [10 min]
        \item Writing Recursive Functions: turning a description of a process into code [15 min]
        \item Recursive Helper Functions: carrying extra state through a helper [12 min]
        \item Tree Recursion: exploring multiple choices with multiple recursive calls [18 min]
        \subitem Don't 
    \end{itemize}
\end{meta}

\section{Recursion Basics}
\begin{questions}
    \subimport{../../topics/recursion/easy/}{three-parts.tex}
    \begin{questionmeta}
        Suggested Time: 3 min; Difficulty: Easy
        \begin{itemize}
            \item Feel free to give an example of a recursive function when explaining this question, and keep getting students to identify parts of recursive functions while solving problems!
        \end{itemize}
    \end{questionmeta}

    \subimport{../../topics/recursion/easy/}{wrong_factorial.tex}
    \begin{questionmeta}
        Suggested Time: 5 min; Difficulty: Easy
        \begin{itemize}
            \item Ask students to trace \lstinline{factorial(3)} to see why it never terminates before asking them to fix it.
            \item Emphasize that the base case's return type should match the type the rest of the function returns.
        \end{itemize}
    \end{questionmeta}

    \subimport{../../topics/recursion/medium/}{sum-prime-digits.tex}
    \begin{questionmeta}
        \begin{itemize}
            \item Again, this problem is likely skippable if you have advanced students.
            \item Make sure your students understand why you are calling \lstinline{sum_prime_digits} twice.
        \end{itemize}    
        Suggested Time: 5 min; Difficulty (Official): Medium; Difficulty (Adjusted): Easy $\rightarrow$ Medium Mezzanine
      \end{questionmeta}
\end{questions}

\newpage
\section{Writing Recursive Functions}
\begin{questions}
    \subimport{../../topics/recursion/easy/}{num-digits.tex}
    \begin{questionmeta}
        Suggested Time: 5 min; Difficulty: Easy
        \begin{itemize}
            \item Good warm-up for describing the recursive case in English first (``the number of digits of $n$ is 1 more than the number of digits of $n$ without its last digit'') before writing code.
        \end{itemize}
    \end{questionmeta}

    \subimport{../../topics/recursion/easy/}{combine.tex}
    \begin{questionmeta}
        Suggested Time: 10 min; Difficulty: Medium
        \begin{itemize}
            \item Walk through \lstinline{combine(43, mul, 2)} concretely before filling in blanks, the same way lecture worked through \lstinline{streak} by picking an example and describing the process in words first.
            \item Common misconception: the order of arguments to \lstinline{f} matters (\lstinline{f(n % 10, result)}, not the reverse).
        \end{itemize}
    \end{questionmeta}
\end{questions}

\section{Recursive Helper Functions}
\begin{questions}
    \subimport{../../topics/recursion/medium/}{positionizer.tex}
    \begin{questionmeta}
        Suggested Time: 10 min (Part a); Part b is a stretch problem, do it only if time allows.
        \begin{itemize}
            \item The given function doesn't have enough information on its own, so we introduce a helper function that takes an extra argument (\lstinline{pos}) to track additional state as we recurse.
            \item Ask: ``What does \lstinline{helper} need to know that \lstinline{positionizer} doesn't pass it directly?'' and ``What should we call \lstinline{helper} with initially?''.
        \end{itemize}
    \end{questionmeta}
\end{questions}

\newpage
\section{Tree Recursion}
% \subimport{../../topics/recursion/text/}{tree_recursion_overview.tex}

\begin{questions}
    \subimport{../../topics/recursion/easy/}{recursion-vs-tree-recursion.tex}
    \begin{questionmeta}
        Suggested time: 3 min
    \end{questionmeta}
    \subimport{../../topics/recursion/easy/}{gibonacci.tex}
    \begin{questionmeta}
        Suggested Time: 5 min; Difficulty: Easy
        \begin{itemize}
            \item This is structurally the same idea as lecture's \lstinline{count_partitions}: at each step we explore using vs.\ not using a given coin size, except here we take a \lstinline{min} over choices instead of summing them.
            \item Ask students what changes if the question instead asked for the \emph{number of ways} to make change rather than the minimum number of coins -- this connects directly back to \lstinline{count_partitions}.
        \end{itemize}
    \end{questionmeta}

    \subimport{../../topics/recursion/medium/}{mario-number.tex}
    \begin{questionmeta}
        Suggested Time: 10 min; Difficulty (Official): Medium;
        The classic mario-number problem from CSM. A worksheet wouldn't be complete without it!
    \end{questionmeta}

    \subimport{../../topics/recursion/medium/}{make-change.tex}
    \begin{questionmeta}
        Suggested Time: 10 min; Difficulty: Medium
        \begin{itemize}
            \item This is structurally the same idea as lecture's \lstinline{count_partitions}: at each step we explore using vs.\ not using a given coin size, except here we take a \lstinline{min} over choices instead of summing them.
            \item Ask students what changes if the question instead asked for the \emph{number of ways} to make change rather than the minimum number of coins -- this connects directly back to \lstinline{count_partitions}.
        \end{itemize}
    \end{questionmeta}

    % \subimport{../../topics/recursion/medium/}{fast_modular_exponentiation.tex}
    % \begin{questionmeta}
    %     Suggested Time: 10 min; Difficulty (Official): Medium; Difficulty (Adjusted): Medium
    % \end{questionmeta}

    % \subimport{../../topics/recursion/medium/}{discussions1.tex}
    % \begin{questionmeta}
    %     Suggested Time: 12 min; Difficulty: Medium
    %     \begin{itemize}
    %         \item Have students identify the two options being explored at each step (an \lstinline{n}-person section or an \lstinline{m}-person section) before writing any code, just like lecture's ``use at least one 4'' vs.\ ``don't use any 4'' split for \lstinline{count_partitions}.
    %         \item Since we only need one option to succeed, the recursive calls are combined with \lstinline{or} rather than added -- a good moment to discuss how the way you combine recursive calls depends on what the function returns.
    %     \end{itemize}
    % \end{questionmeta}
\end{questions}

\end{document}
